\documentclass[12pt,a4paper]{article}
\usepackage{amsmath,amssymb}
\usepackage{graphicx}
\usepackage{listings}
\usepackage{xcolor}
\usepackage{geometry}
\geometry{margin=2.5cm}

% --- تنظیمات نمایش کدها ---
\lstset{
    basicstyle=\small\ttfamily,
    backgroundcolor=\color{gray!10},
    frame=single,
    breaklines=true,
    postbreak=\mbox{\textcolor{red}{$\hookrightarrow$}\space},
    keywordstyle=\color{blue}\bfseries,
    commentstyle=\color{green!60!black},
    stringstyle=\color{purple},
    showstringspaces=false,
    columns=fullflexible,
    keepspaces=true
}

% --- تنظیمات هایپرلینک ---
\usepackage{hyperref}
\hypersetup{
    colorlinks=true,
    linkcolor=blue,
    urlcolor=cyan
}

% --- تنظیمات زی‌پرشین (باید آخرین پکیج باشد) ---
\usepackage{xepersian}
\settextfont{Amiri}
\setlatintextfont{Noto Sans}

\title{\textbf{راهنمای کاربری اجرای سیستم حسگری وای‌فای \lr{(WiFi Sensing Quick Start)}}}
\author{مستندات عملیاتی پلتفرم RuView}
\date{\today}

\begin{document}

\maketitle

\section{مقدمه و پیش‌نیازها}
این راهنما جهت راه‌اندازی سریع و بدون دردسر سیستم حسگری وای‌فای برای کاربران طراحی شده است. فرض بر این است که گره‌های سخت‌افزاری \lr{ESP32-S3} قبلاً با شناسه‌های یکتا (\lr{Node ID}) و \lr{IP} مشخص پروگرام شده‌اند (در صورت نیاز به پروگرام یا تغییر \lr{IP}، به \textbf{گزارش جامع راه‌اندازی سخت‌افزار} مراجعه شود).

\section{گام اول: چیدمان فیزیکی و اتصال گره‌ها به برق}
برای دستیابی به بالاترین دقت در پوشش فضایی و ایجاد شبکه دید متقاطع (\lr{Cross-Beam}):
\begin{enumerate}
    \item ۳ گره \lr{ESP32} را در \textbf{۳ گوشهٔ مختلف اتاق} با حداکثر فاصله ممکن از یکدیگر قرار دهید.
    \item ارتفاع مناسب برای نصب گره‌ها حدود \textbf{۱ تا ۱.۵ متر} (هم‌سطح با بالاتنه افراد) است. از قرار دادن گره‌ها مستقیماً بر روی سطوح فلزی بزرگ خودداری نمایید.
    \item هر گره را با استفاده از کابل \lr{USB} به یک شارژر دیواری، پاوربانک یا درگاه رایانه متصل کنید تا برق آن تأمین شود.
    \item \textbf{یک بار دکمه سخت‌افزاری \lr{RST / EN} را روی هر سه برد فشار دهید} تا حافظه موقت ریست شده و اتصال اولیه به شبکه وای‌فای آغاز گردد.
\end{enumerate}

\section{گام دوم: آماده‌سازی شبکه و فایروال سیستم میزبان}
پیش از اجرای سرور، اطمینان حاصل کنید که رایانه شما به همان شبکه وای‌فای محیط متصل است (یا کابل شبکه متصل است) و پورت دریافت پکت‌های \lr{UDP} روی سیستم‌عامل باز است:

\begin{latin}
\begin{lstlisting}[language=bash]
sudo ufw allow 5005/udp
\end{lstlisting}
\end{latin}

\section{گام سوم: اجرای سرور پردازش \lr{(Sensing Server)}}
یک ترمینال جدید باز کرده و دستور زیر را برای اجرای سرور پردازش به همراه تنظیمات اسلات‌های زمانی (\lr{TDM}) جهت هماهنگی ۳ گره اجرا نمایید:

\begin{latin}
\begin{lstlisting}[language=bash]
cd ~/Desktop/Amirali/WifiSensing/ruview/v2
WDP_TDM_SLOTS=3 WDP_TDM_SLOT_US=200000 cargo run --release --package wifi-densepose-sensing-server -- \
  --udp-bind 0.0.0.0 \
  --udp-insecure-lan \
  --source esp32
\end{lstlisting}
\end{latin}

\paragraph{نکته:} در صورت تمایل به فیلتر کردن تردد افراد در راهرو یا بیرون اتاق، پارامتر \lr{\texttt{WDP\_MOTION\_THRESHOLD=1.5}} را نیز به ابتدای دستور فوق اضافه کنید.

\section{گام چهارم: مشاهده خروجی در داشبورد وب}
پس از بالا آمدن سرور در ترمینال، مرورگر خود (مانند فایرفاکس یا کروم) را باز کرده و آدرس زیر را وارد کنید:

\begin{latin}
\begin{center}
\textbf{\url{http://localhost:8080/ui/index.html}}
\end{center}
\end{latin}

\subsection{بررسی تب‌های اصلی داشبورد:}
\begin{itemize}
    \item \textbf{تب \lr{Sensing}:} بررسی وضعیت اتصال گره‌ها؛ مطمئن شوید هر ۳ گره شناسایی شده و دربرای برسی node های متصل میتوانید در این بخش به 
    \lr{Node Status} بروید.
    \item \textbf{تب \lr{Sensing}:} مشاهده جریان زنده تغییرات سیگنال، سطح \lr{RSSI}، میزان واریانس حرکت افراد و نمودار تنفس لحظه‌ای.
\end{itemize}

اکنون سیستم کاملاً آماده استفاده و پایش محیط است.

\end{document}